\chapter{Dater à Rome au IV\textsuperscript{e}~siècle}\label{ch:dater}

\section{Consulats et post-consulats}

La datation absolue des épitaphes du corpus repose presque exclusivement sur la mention des consuls ordinaires. Deux difficultés sont bien connues. Le nom des consuls de l’année pouvait n’être pas encore connu, ou pas reconnu, à Rome : on datait alors « après le consulat » des consuls précédents (\emph{post consulatum}). D’autre part, la proclamation des consuls dépendait de la situation politique : usurpations, consulats non reconnus dans une partie de l’Empire, condamnation de la mémoire d’un consul. L’instrument de référence pour ces questions est l’ouvrage de Bagnall, Cameron, Schwartz et Worp ; il est cité ici à travers les fiches de l’EDB lorsque celles-ci le reproduisent.

Le corpus illustre plusieurs de ces situations, et ses textes permettent d’en suivre les effets.

\emph{346.} Les quatre épitaphes classées en 346 datent par le post-consulat d’Amantius et Albinus (\nnums{15}–18). Le \nnum{15} associe une naissance datée par le consulat de Sabinus et Rufinus (316) et un décès « après le consulat d’Amantius et d’Albinus » : le calcul de la durée de vie confirme que ce post-consulat désigne bien 346 (chapitre~\ref{ch:calendrier}).

\emph{352–353.} Les consuls désignés sous l’usurpation de Magnence apparaissent en 352 (Décence et Paulus : \nnums{12}, 37–39) et, selon la lecture retenue, en 353 (Magnence et Décence : \nnum{40}). L’usage romain des consuls de Magnence est relevé par la recherche récente, qui note que la plupart des exemples datés de 352 remontent au printemps. Au \nnum{40}, la formule \emph{M⸢a⸣nnentio et Decen[ti]o} est interprétée par l’EDB, d’après la CLRE, comme une formule de 352 mal combinée par le lapicide ; la datation est déclarée contestée (chapitre~\ref{ch:fondement}). Le \nnum{12} porte deux formules consulaires : \emph{Amantio et Albino} (345), qui date un événement antérieur non précisé, et \emph{[Decentio et P]aulo} (352), rattachée à la déposition ; seule la fin du nom de Paulus est lue.

\emph{363.} Les dix épitaphes de 363 (\nnums{73}–82) datent par le quatrième consulat de Julien, avec Sallustius. Plusieurs sont postérieures à la mort de l’empereur, survenue en juin : le \nnum{81} est daté de décembre. Le consulat datait l’année, quel que fût le sort de l’empereur.

\emph{375.} Les dix textes du corpus classés en 375 (\nnums{179}–188) datent par le post-consulat de Gratien III et d’Equitius, ce qui concorde avec l’absence de nouveau collège dans la tradition des fastes (point à contrôler sur la CLRE). Au \nnum{180}, la formule n’est transmise que par l’EDB.

\emph{388.} Les textes datent soit par le consulat de l’usurpateur Magnus Maximus (\nnums{295}–297, 300, 302 et 311), soit par Théodose et Merobaudes. Au \nnum{298}, Merobaudes est dit consul pour la troisième fois (\emph{Merobaude v(iro) c(larissimo) III}, lu) ; au \nnum{303}, pour la deuxième. L’existence d’un troisième consulat de Merobaudes en 388 pose un problème prosopographique que la PLRE et la CLRE discutent, et qui n’a pas pu être contrôlé ici.

\emph{393.} Les épitaphes datent par Théodose et l’usurpateur Eugène (\nnums{351}–361). Le \nnum{351} porte en toutes lettres \emph{Theodosio III et Eugenio Augg(ustis)} ; le \nnum{360} transpose la formule en grec (\emph{Θεοδοσίῳ τό […]}).

\emph{394.} Le corpus atteste à la fois le consulat de Virius Nicomachus Flavianus, reconnu seulement dans le territoire d’Eugène\footnote{J. J. O’Donnell, « The Career of Virius Nicomachus Flavianus », \emph{Phoenix}, 32/2, 1978, p. 129–143, DOI 10.2307/1087470 (métadonnées vérifiées). Au \nnum{362}, le nom est en grande partie restitué : seul \emph{-abiani(!) v(iri)} est lu.} (\nnums{362}, 364–366), et le collège oriental d’Arcadius III et Honorius II (\nnums{363} et 367).

\emph{399.} Seul Mallius Theodorus est nommé (\nnums{429}, 430, 432, 433, 436–438), le consulat oriental de l’eunuque Eutrope n’ayant pas été reconnu en Occident. Au \nnum{435}, le nom n’est qu’une restitution de l’EDB ; au \nnum{431}, la formule \emph{post consulatu … Arcadi IIII et Eutychiani} (avec \emph{Harchadi} gravé pour \emph{Honorii}) est attribuée par la CLRE à 398.

\emph{400.} Stilichon est seul nommé (\nnums{439}–446). Au \nnum{444}, le groupe \emph{[iter(um) con]sule} est restitué par l’éditeur ; s’il était exact, il désignerait le second consulat de Stilichon (405). L’année 400 est conservée sous réserve. Au \nnum{445}, l’EDB restitue \emph{[iterum?]} avec doute.

\section{Table des consulats}

La table~\ref{tab:fastes} donne, pour chaque année, le collège consulaire selon la tradition des fastes et le nombre de textes de l’ensemble B classés sous cette année. Elle a servi à l’identification des consuls. Elle n’a pas été collationnée sur la CLRE et doit l’être avant publication ; les points litigieux sont signalés au chapitre suivant.

% généré par scripts/04_generated.py
\begin{longtable}{@{}r p{0.70\linewidth} r@{}}
\caption{Collèges consulaires 340–403 selon la tradition des fastes et nombre de textes distincts (ensemble B) classés sous chaque année. Table non collationnée sur la CLRE.}\label{tab:fastes}\\
\toprule Année & Collège consulaire & Textes\\\midrule\endfirsthead
\toprule Année & Collège consulaire & Textes\\\midrule\endhead
\bottomrule\endfoot
340 & Acindynus et Proculus & 3 \\
341 & Marcellinus et Probinus & 1 \\
342 & Constantius Aug. III et Constans Aug. II & 2 \\
343 & Placidus et Romulus & 3 \\
344 & Leontius et Sallustius & 1 \\
345 & Amantius et Albinus & 3 \\
346 & Constantius Aug. IV et Constans Aug. III (à Rome : p.c. Amantii et Albini) & 4 \\
347 & Rufinus et Eusebius & 3 \\
348 & Philippus et Salia & 8 \\
349 & Limenius et Catullinus & 4 \\
350 & Sergius et Nigrinianus & 2 \\
351 & Occident (Magnence) : Magnentius Aug. et Gaiso ; Orient : p.c. Sergii et Nigriniani & 1 \\
352 & Occident (Magnence) : Decentius Caes. et Paulus ; Constance : Constantius Aug. V et Constantius Caes. & 4 \\
353 & Occident (Magnence) : Magnentius Aug. II et Decentius Caes. II ; Constance : Constantius Aug. VI et Constantius Caes. II & 1 \\
354 & Constantius Aug. VII et Constantius Caes. III & 0 \\
355 & Arbitio et Lollianus & 4 \\
356 & Constantius Aug. VIII et Iulianus Caes. & 2 \\
357 & Constantius Aug. IX et Iulianus Caes. II & 2 \\
358 & Datianus et Cerealis & 5 \\
359 & Eusebius et Hypatius & 2 \\
360 & Constantius Aug. X et Iulianus Caes. III & 3 \\
361 & Taurus et Florentius & 3 \\
362 & Mamertinus et Nevitta & 11 \\
363 & Iulianus Aug. IV et Sallustius & 10 \\
364 & Iovianus Aug. et Varronianus & 9 \\
365 & Valentinianus Aug. et Valens Aug. & 6 \\
366 & Gratianus nob. puer et Dagalaifus & 7 \\
367 & Lupicinus et Iovinus & 18 \\
368 & Valentinianus Aug. II et Valens Aug. II & 7 \\
369 & Valentinianus nob. puer et Victor & 4 \\
370 & Valentinianus Aug. III et Valens Aug. III & 8 \\
371 & Gratianus Aug. II et Probus & 12 \\
372 & Modestus et Arintheus & 11 \\
373 & Valentinianus Aug. IV et Valens Aug. IV & 6 \\
374 & Gratianus Aug. III et Equitius & 7 \\
375 & Pas de consuls ordinaires : p.c. Gratiani III et Equitii & 10 \\
376 & Valens Aug. V et Valentinianus iun. Aug. & 6 \\
377 & Gratianus Aug. IV et Merobaudes & 11 \\
378 & Valens Aug. VI et Valentinianus iun. Aug. II & 3 \\
379 & Ausonius et Olybrius & 3 \\
380 & Gratianus Aug. V et Theodosius Aug. & 6 \\
381 & Syagrius et Eucherius & 10 \\
382 & Antonius et Syagrius & 13 \\
383 & Merobaudes II et Saturninus & 20 \\
384 & Richomeres et Clearchus & 15 \\
385 & Arcadius Aug. et Bauto & 6 \\
386 & Honorius nob. puer et Euodius & 8 \\
387 & Valentinianus Aug. III et Eutropius & 4 \\
388 & Theodosius Aug. II et Cynegius (Occident : Magnus Maximus Aug. II) ; Merobaudes III attesté par l'épigraphie (n° 298) & 10 \\
389 & Timasius et Promotus & 13 \\
390 & Valentinianus Aug. IV et Neoterius & 10 \\
391 & Tatianus et Symmachus & 13 \\
392 & Arcadius Aug. II et Rufinus & 8 \\
393 & Theodosius Aug. III et Abundantius ; à Rome : Eugenius Aug. & 11 \\
394 & Arcadius Aug. III et Honorius Aug. II ; territoire d'Eugène : Nicomachus Flavianus & 6 \\
395 & Olybrius et Probinus & 17 \\
396 & Arcadius Aug. IV et Honorius Aug. III & 10 \\
397 & Caesarius et Atticus & 22 \\
398 & Honorius Aug. IV et Eutychianus & 12 \\
399 & Mallius Theodorus et Eutropius (Occident : Theodorus seul) & 9 \\
400 & Stilicho et Aurelianus & 8 \\
401 & Vincentius et Fravitta & 13 \\
402 & Arcadius Aug. V et Honorius Aug. V & 9 \\
403 & Theodosius iun. Aug. et Rumoridus & 10 \\
\end{longtable}


\section{Le calendrier romain et la conversion des dates}

Les quantièmes sont exprimés selon le calendrier romain, par décompte inclusif à partir des calendes (le 1\textsuperscript{er} du mois), des nones (le 5, ou le 7 en mars, mai, juillet et octobre) et des ides (le 13, ou le 15 les mêmes mois). « Le 15\textsuperscript{e} jour avant les calendes d’avril » désigne le 18 mars. En année bissextile julienne, le 6\textsuperscript{e} jour avant les calendes de mars est redoublé ; la conversion en tient compte. Les dates sont converties dans le calendrier julien.

\begin{table}[tbp]\centering\small
\caption{Le calendrier romain (année ordinaire). Décompte inclusif : \emph{a.d.} N \emph{Kal.} = premier jour du mois suivant moins (N − 1) jours.}\label{tab:calendrier}
\begin{tabular}{@{}l r r r l@{}}\toprule Mois & Kal. & Non. & Id. & Jours suivant les ides\\\midrule
Janvier & 1 & 5 & 13 & 14 = \emph{a.d.} XIX \emph{Kal. Febr.} … 31 = \emph{prid. Kal. Febr.}\\
Février & 1 & 5 & 13 & 14 = \emph{a.d.} XVI \emph{Kal. Mart.} … 28 = \emph{prid. Kal. Mart.}\\
Mars & 1 & 7 & 15 & 16 = \emph{a.d.} XVII \emph{Kal. Apr.} … 31 = \emph{prid. Kal. Apr.}\\
Avril & 1 & 5 & 13 & 14 = \emph{a.d.} XVIII \emph{Kal. Mai.} … 30 = \emph{prid. Kal. Mai.}\\
Mai & 1 & 7 & 15 & 16 = \emph{a.d.} XVII \emph{Kal. Iun.} … 31 = \emph{prid. Kal. Iun.}\\
Juin & 1 & 5 & 13 & 14 = \emph{a.d.} XVIII \emph{Kal. Iul.} … 30 = \emph{prid. Kal. Iul.}\\
Juillet & 1 & 7 & 15 & 16 = \emph{a.d.} XVII \emph{Kal. Aug.} … 31 = \emph{prid. Kal. Aug.}\\
Août & 1 & 5 & 13 & 14 = \emph{a.d.} XIX \emph{Kal. Sept.} … 31 = \emph{prid. Kal. Sept.}\\
Septembre & 1 & 5 & 13 & 14 = \emph{a.d.} XVIII \emph{Kal. Oct.} … 30 = \emph{prid. Kal. Oct.}\\
Octobre & 1 & 7 & 15 & 16 = \emph{a.d.} XVII \emph{Kal. Nov.} … 31 = \emph{prid. Kal. Nov.}\\
Novembre & 1 & 5 & 13 & 14 = \emph{a.d.} XVIII \emph{Kal. Dec.} … 30 = \emph{prid. Kal. Dec.}\\
Décembre & 1 & 5 & 13 & 14 = \emph{a.d.} XIX \emph{Kal. Ian.} … 31 = \emph{prid. Kal. Ian.}\\\bottomrule
\end{tabular}
\end{table}

\section{Jours de la semaine et âge de la lune}

Quelques textes indiquent le jour de la semaine selon la nomenclature planétaire (\emph{dies Solis, Lunae, Martis, Mercurii, Iovis, Veneris, Saturni}) et, plus rarement, l’âge de la lune. Le jour de la semaine est vérifié par le calcul du jour julien dans le calendrier julien. L’âge de la lune et sa position dans le zodiaque sont calculés avec la bibliothèque d’éphémérides PyEphem (version 4.2.1), dont on a vérifié qu’elle emploie le calendrier julien pour ces dates : le jour julien calculé pour le 23 août 386 coïncide avec celui de l’algorithme de conversion, soit 1\,862\,279. Ces contrôles font l’objet du chapitre~\ref{ch:calendrier}.
